\documentclass[10pt,a4paper]{amsart}
\usepackage{amsmath,amssymb,amsthm,graphicx}
% Retained preamble for the Stacks Project chapter snapshot in this directory.
% The chapter's own preamble.tex is a copy of the shared Stacks Project preamble; the block
% below is sliced MECHANICALLY, line for line, from it. Nothing was added, removed or
% redefined; the wrapper class, the xr-hyper/hyperref wiring and the externaldocument list
% are omitted because the assembler supplies the document class.
%
% lines taken from preamble.tex: 3, 5, 9, 11, 12, 156, 157, 158, ... (51 lines total)
%
% package coverage restored from the chapter's own preamble (2026-10-13): xy, xy-2cell

\usepackage[all]{xy}
\xyoption{2cell}
\UseAllTwocells

\newenvironment{reference}{\comment}{\endcomment}
\newenvironment{slogan}{\comment}{\endcomment}
\newenvironment{history}{\comment}{\endcomment}
\theoremstyle{plain}
\newtheorem{theorem}[subsection]{Theorem}
\newtheorem{proposition}[subsection]{Proposition}
\newtheorem{lemma}[subsection]{Lemma}
\theoremstyle{definition}
\newtheorem{definition}[subsection]{Definition}
\newtheorem{example}[subsection]{Example}
\newtheorem{exercise}[subsection]{Exercise}
\newtheorem{situation}[subsection]{Situation}
\theoremstyle{remark}
\newtheorem{remark}[subsection]{Remark}
\newtheorem{remarks}[subsection]{Remarks}
\numberwithin{equation}{subsection}
\def\lim{\mathop{\mathrm{lim}}\nolimits}
\def\colim{\mathop{\mathrm{colim}}\nolimits}
\def\Spec{\mathop{\mathrm{Spec}}}
\def\Hom{\mathop{\mathrm{Hom}}\nolimits}
\def\Ext{\mathop{\mathrm{Ext}}\nolimits}
\def\SheafHom{\mathop{\mathcal{H}\!\mathit{om}}\nolimits}
\def\SheafExt{\mathop{\mathcal{E}\!\mathit{xt}}\nolimits}
\def\Sch{\mathit{Sch}}
\def\Mor{\mathop{\mathrm{Mor}}\nolimits}
\def\Ob{\mathop{\mathrm{Ob}}\nolimits}
\def\Sh{\mathop{\mathit{Sh}}\nolimits}
\def\NL{\mathop{N\!L}\nolimits}
\def\CH{\mathop{\mathrm{CH}}\nolimits}
\def\proetale{{pro\text{-}\acute{e}tale}}
\def\etale{{\acute{e}tale}}
\def\QCoh{\mathit{QCoh}}
\def\Ker{\mathop{\mathrm{Ker}}}
\def\Im{\mathop{\mathrm{Im}}}
\def\Coker{\mathop{\mathrm{Coker}}}
\def\Coim{\mathop{\mathrm{Coim}}}
\DeclareMathSymbol{\boxtimes}{\mathbin}{AMSa}{"02}
\def\QCohstack{\mathcal{QC}\!\mathit{oh}}
\def\Cohstack{\mathcal{C}\!\mathit{oh}}
\def\Spacesstack{\mathcal{S}\!\mathit{paces}}
\def\Quotfunctor{\mathrm{Quot}}
\def\Hilbfunctor{\mathrm{Hilb}}
\def\Curvesstack{\mathcal{C}\!\mathit{urves}}
\def\Polarizedstack{\mathcal{P}\!\mathit{olarized}}
\def\Complexesstack{\mathcal{C}\!\mathit{omplexes}}
\def\Pic{\mathop{\mathrm{Pic}}\nolimits}
\def\Picardstack{\mathcal{P}\!\mathit{ic}}
\def\Picardfunctor{\mathrm{Pic}}
\def\Deformationcategory{\mathcal{D}\!\mathit{ef}}
\usepackage{cleveref}
\crefname{originallemma}{original lemma}{original lemmas}
\Crefname{originallemma}{Original lemma}{Original lemmas}
\makeatletter
\newlabel{algebra-lemma-annihilator-flat-base-change}{{07T8}{source}{}{sourcepointer-1}{}}
\newlabel{algebra-lemma-annihilator-flat-base-change@cref}{{[originallemma][0][]07T8}{[1][0][]source}{}{}{}}
\makeatother
\crefname{originallemma}{original lemma}{original lemmas}
\Crefname{originallemma}{Original lemma}{Original lemmas}
\makeatletter
\newlabel{algebra-lemma-support-base-change}{{0BUR}{source}{}{sourcepointer-2}{}}
\newlabel{algebra-lemma-support-base-change@cref}{{[originallemma][0][]0BUR}{[1][0][]source}{}{}{}}
\makeatother
\crefname{originallemma}{original lemma}{original lemmas}
\Crefname{originallemma}{Original lemma}{Original lemmas}
\makeatletter
\newlabel{algebra-lemma-support-closed}{{00L2}{source}{}{sourcepointer-3}{}}
\newlabel{algebra-lemma-support-closed@cref}{{[originallemma][0][]00L2}{[1][0][]source}{}{}{}}
\makeatother
\crefname{originallemma}{original lemma}{original lemmas}
\Crefname{originallemma}{Original lemma}{Original lemmas}
\makeatletter
\newlabel{lemma-dualizing-base-change-of-flat}{{0BVF}{source}{}{sourcepointer-4}{}}
\newlabel{lemma-dualizing-base-change-of-flat@cref}{{[originallemma][0][]0BVF}{[1][0][]source}{}{}{}}
\makeatother
\crefname{originallemma}{original lemma}{original lemmas}
\Crefname{originallemma}{Original lemma}{Original lemmas}
\makeatletter
\newlabel{lemma-trace-base-change}{{0BT9}{source}{}{sourcepointer-5}{}}
\newlabel{lemma-trace-base-change@cref}{{[originallemma][0][]0BT9}{[1][0][]source}{}{}{}}
\makeatother
\title{Stacks, Discriminants: lemma-base-change-different}
\author{The Stacks Project Authors}
\date{stacks-project:discriminant.tex:local-snapshot-20261012; GNU FDL 1.2 or later}
\begin{document}
\maketitle
\noindent\textit{Editorial note.} Statement span is lines 1812-1827 of sources/stacks-discriminant/pinned-original-source.tex; proof span is lines 1829-1860, the complete author proof (30 body lines).\par\medskip
\noindent\textit{Editorial note.} Navigation stubs: 5 cited-but-not-included labels are supplied as non-load-bearing pointers with their pinned tags where the snapshot tag index has them.\par\medskip
\noindent\textit{Editorial note.} extra\_packages is empty on purpose: the retained preamble already loads xy with options and a second declaration clashes in the assembler wrapper.\par\medskip
\noindent\textit{Editorial note.} Text is retained verbatim, including the author's own typos and placeholders.\par\medskip
\noindent\textit{Editorial note.} Queue index 15 of queue\_discriminant.json (pb=1456, stmt 1812-1827, proof 1829-1860).\par\medskip
\noindent\textit{Allowed standard prerequisites (editorial).}\begin{itemize}
\item Explicit prerequisite boundary. The cited inputs are prerequisites; the argument in the retained span is not replaced by any cited result.
\end{itemize}\medskip
\ifdefined\hypertarget\hypertarget{sourcepointer-1}{}\fi
\noindent\textit{Source reference.} Original lemma 07T8: sources/stacks-discriminant/pinned-original-source.tex, source lines 1860--1860. Cited by the retained statement or proof but not part of the collected spans; supplied only as a navigation pointer. The argument does not depend on its content.\par\medskip
\ifdefined\hypertarget\hypertarget{sourcepointer-2}{}\fi
\noindent\textit{Source reference.} Original lemma 0BUR: sources/stacks-discriminant/pinned-original-source.tex, source lines 1860--1860. Cited by the retained statement or proof but not part of the collected spans; supplied only as a navigation pointer. The argument does not depend on its content.\par\medskip
\ifdefined\hypertarget\hypertarget{sourcepointer-3}{}\fi
\noindent\textit{Source reference.} Original lemma 00L2: sources/stacks-discriminant/pinned-original-source.tex, source lines 1860--1860. Cited by the retained statement or proof but not part of the collected spans; supplied only as a navigation pointer. The argument does not depend on its content.\par\medskip
\ifdefined\hypertarget\hypertarget{sourcepointer-4}{}\fi
\noindent\textit{Source reference.} Original lemma 0BVF: sources/stacks-discriminant/pinned-original-source.tex, source lines 590--593. Cited by the retained statement or proof but not part of the collected spans; supplied only as a navigation pointer. The argument does not depend on its content.\par\medskip
\ifdefined\hypertarget\hypertarget{sourcepointer-5}{}\fi
\noindent\textit{Source reference.} Original lemma 0BT9: sources/stacks-discriminant/pinned-original-source.tex, source lines 894--907. Cited by the retained statement or proof but not part of the collected spans; supplied only as a navigation pointer. The argument does not depend on its content.\par\medskip
\begin{lemma}
\label{lemma-base-change-different}
Consider a cartesian diagram of Noetherian schemes
$$
\xymatrix{
Y' \ar[d]_{f'} \ar[r] & Y \ar[d]^f \\
X' \ar[r]^g & X
}
$$
with $f$ flat and quasi-finite. Let $R \subset Y$, resp.\ $R' \subset Y'$
be the closed subscheme cut out by the different
$\mathfrak{D}_f$, resp.\ $\mathfrak{D}_{f'}$.
Then $Y' \to Y$ induces a bijective closed immersion $R' \to R \times_Y Y'$.
If $g$ is flat or if $\omega_{Y/X}$ is invertible, then
$R' = R \times_Y Y'$.
\end{lemma}

\begin{proof}
There is an immediate reduction to the case where $X$, $X'$, $Y$, $Y'$
are affine. In other words, we have a cocartesian diagram of Noetherian
rings
$$
\xymatrix{
B' & B \ar[l] \\
A' \ar[u] & A \ar[l] \ar[u]
}
$$
with $A \to B$ flat and quasi-finite. The base change map
$\omega_{B/A} \otimes_B B' \to \omega_{B'/A'}$ is an isomorphism
(Lemma \ref{lemma-dualizing-base-change-of-flat}) and maps
the trace element $\tau_{B/A}$ to the trace element $\tau_{B'/A'}$
(Lemma \ref{lemma-trace-base-change}).
Hence the finite $B$-module $Q = \Coker(\tau_{B/A} : B \to \omega_{B/A})$
satisfies $Q \otimes_B B' = \Coker(\tau_{B'/A'} : B' \to \omega_{B'/A'})$.
Thus $\mathfrak{D}_{B/A}B' \subset \mathfrak{D}_{B'/A'}$ which means
we obtain the closed immersion $R' \to R \times_Y Y'$.
Since $R = \text{Supp}(Q)$ and $R' = \text{Supp}(Q \otimes_B B')$
(Algebra, Lemma \ref{algebra-lemma-support-closed})
we see that $R' \to R \times_Y Y'$ is bijective by
Algebra, Lemma \ref{algebra-lemma-support-base-change}.
The equality $\mathfrak{D}_{B/A}B' = \mathfrak{D}_{B'/A'}$ holds
if $B \to B'$ is flat, e.g., if $A \to A'$ is flat, see
Algebra, Lemma \ref{algebra-lemma-annihilator-flat-base-change}.
Finally, if $\omega_{B/A}$ is invertible, then we can localize
and assume $\omega_{B/A} = B \lambda$. Writing $\tau_{B/A} = b\lambda$
we see that $Q = B/bB$ and $\mathfrak{D}_{B/A} = bB$.
The same reasoning over $B'$
gives $\mathfrak{D}_{B'/A'} = bB'$ and the lemma is proved.
\end{proof}

\end{document}
